% ============================================================
\chapter{Courbure --- Tenseur de Riemann}
\label{ch:courbure}
% ============================================================

La courbure est la quantit\'e qui distingue une g\'eom\'etrie d'une autre. Sur un plan, les droites parall\`eles restent parall\`eles~; sur une sph\`ere, elles finissent par se rencontrer~; sur une surface en selle de cheval, elles divergent. Riemann a compris, d\`es 1854, que cette d\'eviation des g\'eod\'esiques est mesur\'ee par un objet tensoriel --- le \emph{tenseur de courbure} --- qui encode, en chaque point et dans chaque plan tangent, comment l'espace se courbe. Ce tenseur, formalis\'e par Christoffel et Ricci-Curbastro, est l'ingr\'edient central de la relativit\'e g\'en\'erale d'Einstein~: la courbure de l'espace-temps \emph{est} la gravitation.

\section{Le tenseur de courbure de Riemann}

\begin{definition}[Tenseur de courbure]
Le \emph{tenseur de courbure de Riemann} de la connexion de Levi-Civita $\nabla$
est l'application $R : \mathfrak{X}(M)^3 \to \mathfrak{X}(M)$ d\'efinie par~:
\[
    R(X,Y)Z = \nabla_X \nabla_Y Z - \nabla_Y \nabla_X Z - \nabla_{[X,Y]} Z.
\]
\end{definition}

\begin{intuition}[title=\textbf{La courbure mesure la non-commutativit\'e}]
Le tenseur de Riemann mesure le d\'efaut de commutativit\'e de la d\'erivation
covariante~: si l'espace \'etait plat, on aurait $\nabla_X \nabla_Y = \nabla_Y \nabla_X$
(modulo le crochet de Lie). La courbure est l'obstruction \`a ce que le transport
parall\`ele soit ind\'ependant du chemin.
\end{intuition}

\begin{proposition}
$R$ est un tenseur de type $(1,3)$~: il est $C^\infty(M)$-multilin\'eaire en $X$,
$Y$ et $Z$.
\end{proposition}

\begin{proof}
V\'erifions la $C^\infty(M)$-lin\'earit\'e en $Z$ (les autres sont similaires)~:
\begin{align*}
    R(X,Y)(fZ) &= \nabla_X\nabla_Y(fZ) - \nabla_Y\nabla_X(fZ) - \nabla_{[X,Y]}(fZ) \\
    &= \nabla_X\big((Yf)Z + f\nabla_Y Z\big) - \nabla_Y\big((Xf)Z + f\nabla_X Z\big)
    - ([X,Y]f)Z - f\nabla_{[X,Y]}Z.
\end{align*}
En d\'eveloppant et simplifiant~: les termes en $Xf$, $Yf$, $XYf$, $YXf$ se
compensent (en utilisant $[X,Y]f = XYf - YXf$), et il reste $fR(X,Y)Z$.
\end{proof}

\section{Expression en coordonn\'ees}

En coordonn\'ees locales, $R(\partial_i, \partial_j)\partial_k = R^l_{kij}\, \partial_l$,
avec~:
\[
    R^l_{kij} = \partial_i \Gamma^l_{jk} - \partial_j \Gamma^l_{ik}
    + \Gamma^l_{im}\Gamma^m_{jk} - \Gamma^l_{jm}\Gamma^m_{ik}.
\]

On d\'efinit \'egalement le \emph{tenseur de courbure abaiss\'e}~:
\[
    R_{ijkl} = g_{im} R^m_{jkl},
\]
de sorte que $R_{ijkl} = \inner{R(\partial_k, \partial_l)\partial_j}{\partial_i}$.

\begin{formules}[title=\textbf{Composantes du tenseur de Riemann}]
\[
    R^l_{kij} = \partial_i \Gamma^l_{jk} - \partial_j \Gamma^l_{ik}
    + \Gamma^l_{im}\Gamma^m_{jk} - \Gamma^l_{jm}\Gamma^m_{ik}.
\]
En coordonn\'ees normales en $p$~: $R^l_{kij}(p) = \partial_i \Gamma^l_{jk}(p)
- \partial_j \Gamma^l_{ik}(p)$.
\end{formules}

\section{Sym\'etries du tenseur de Riemann}

\begin{theoreme}[Sym\'etries fondamentales]
\label{thm:symmetries}
Le tenseur de courbure abaiss\'e $R_{ijkl}$ poss\`ede les sym\'etries suivantes~:
\begin{enumerate}
    \item \textbf{Antisym\'etrie en $(k,l)$~:} $R_{ijkl} = -R_{ijlk}$.
    \item \textbf{Antisym\'etrie en $(i,j)$~:} $R_{ijkl} = -R_{jikl}$.
    \item \textbf{Sym\'etrie par paires~:} $R_{ijkl} = R_{klij}$.
    \item \textbf{Premi\`ere identit\'e de Bianchi (alg\'ebrique)~:}
          $R_{ijkl} + R_{iklj} + R_{iljk} = 0$.
\end{enumerate}
En notation intrins\`eque~:
\begin{enumerate}
    \item $R(X,Y) = -R(Y,X)$.
    \item $\inner{R(X,Y)Z}{W} = -\inner{R(X,Y)W}{Z}$.
    \item $\inner{R(X,Y)Z}{W} = \inner{R(Z,W)X}{Y}$.
    \item $R(X,Y)Z + R(Y,Z)X + R(Z,X)Y = 0$.
\end{enumerate}
\end{theoreme}

\begin{proof}
(1) est imm\'ediat par d\'efinition ($R(X,Y) = -R(Y,X)$).

(2) D\'ecoule de la compatibilit\'e avec la m\'etrique. En effet~:
\begin{align*}
    \inner{R(X,Y)Z}{Z} &= \inner{\nabla_X\nabla_Y Z - \nabla_Y\nabla_X Z}{Z} \\
    &= X\inner{\nabla_Y Z}{Z} - \inner{\nabla_Y Z}{\nabla_X Z}
    - Y\inner{\nabla_X Z}{Z} + \inner{\nabla_X Z}{\nabla_Y Z} \\
    &= \frac{1}{2}XY\inner{Z}{Z} - \frac{1}{2}YX\inner{Z}{Z}
    = \frac{1}{2}[X,Y]\inner{Z}{Z} = \inner{\nabla_{[X,Y]}Z}{Z}.
\end{align*}
Donc $\inner{R(X,Y)Z}{Z} = 0$, ce qui par polarisation donne (2).

(3) d\'ecoule de (1), (2) et (4).

(4) Premi\`ere identit\'e de Bianchi~: en utilisant l'absence de torsion et la
d\'efinition, on d\'eveloppe les trois termes cycliques et on constate que la somme
s'annule.
\end{proof}

\begin{proposition}[Nombre de composantes ind\'ependantes]
En dimension $n$, le nombre de composantes ind\'ependantes du tenseur de Riemann est~:
\[
    \frac{n^2(n^2-1)}{12}.
\]
Pour $n = 2$~: $1$, pour $n = 3$~: $6$, pour $n = 4$~: $20$.
\end{proposition}

\section{Identit\'es de Bianchi}

\begin{theoreme}[Seconde identit\'e de Bianchi (diff\'erentielle)]
\[
    (\nabla_X R)(Y,Z) + (\nabla_Y R)(Z,X) + (\nabla_Z R)(X,Y) = 0.
\]
En coordonn\'ees~: $\nabla_m R^l_{kij} + \nabla_i R^l_{kmj} + \nabla_j R^l_{kim} = 0$.
\end{theoreme}

\begin{proof}
On travaille en coordonn\'ees normales au point $p$ o\`u $\Gamma^k_{ij}(p) = 0$.
La d\'eriv\'ee covariante de $R$ se r\'eduit \`a la d\'eriv\'ee partielle, et le r\'esultat
d\'ecoule d'un calcul direct utilisant la commutation des d\'eriv\'ees partielles.
\end{proof}

\begin{remarque}
La seconde identit\'e de Bianchi joue un r\^ole fondamental~: contract\'ee, elle
donne la conservation du tenseur d'Einstein $G_{ij} = R_{ij} - \frac{1}{2}Sg_{ij}$
en relativit\'e g\'en\'erale.
\end{remarque}

\section{Courbure et transport parall\`ele}

\begin{theoreme}[Interpr\'etation via le transport parall\`ele]
Soient $X, Y$ deux champs de vecteurs qui commutent ($[X,Y] = 0$). Consid\'erons
le parall\'elogramme infinit\'esimal form\'e par les flots de $sX$ et $tY$. Le
transport parall\`ele d'un vecteur $Z$ le long de ce parall\'elogramme donne, au
second ordre~:
\[
    Z_{\text{final}} - Z_{\text{initial}} = st\, R(X,Y)Z + O(s^2 + t^2).
\]
\end{theoreme}

\begin{corollaire}[Platitude]
$(M, g)$ est plate ($R \equiv 0$) si et seulement si le transport parall\`ele
est ind\'ependant du chemin, si et seulement si $(M, g)$ est localement isom\'etrique
\`a $\R^n$.
\end{corollaire}

\section{Champs de Jacobi}

\begin{definition}[Champ de Jacobi]
Soit $\gamma$ une g\'eod\'esique. Un champ de vecteurs $J$ le long de $\gamma$ est
un \emph{champ de Jacobi} s'il satisfait l'\emph{\'equation de Jacobi}~:
\[
    \frac{D^2 J}{dt^2} + R(\dot\gamma, J)\dot\gamma = 0.
\]
\end{definition}

\begin{proposition}
Les champs de Jacobi le long d'une g\'eod\'esique forment un espace vectoriel de
dimension $2n$, d\'etermin\'e par les conditions initiales $J(0)$ et $\frac{DJ}{dt}(0)$.
\end{proposition}

\begin{theoreme}[Origine des champs de Jacobi]
Les champs de Jacobi sont les champs de variation de familles de g\'eod\'esiques~:
si $\gamma_s$ est une famille lisse de g\'eod\'esiques, alors
$J(t) = \frac{\partial\gamma_s}{\partial s}\big|_{s=0}$ est un champ de Jacobi
le long de $\gamma_0$.
\end{theoreme}

\begin{proof}
Posons $T = \dot\gamma_s$ et $J = \frac{\partial\gamma_s}{\partial s}$.
Comme $[T, J] = 0$ (sym\'etrie de la variation) et $\nabla_T T = 0$ (\'equation
des g\'eod\'esiques)~:
\[
    \frac{D^2 J}{dt^2} = \nabla_T\nabla_T J = \nabla_T\nabla_J T
    = \nabla_J\nabla_T T + R(T,J)T = R(T,J)T = -R(T,J)T.
\]
D'o\`u $\frac{D^2 J}{dt^2} + R(T, J)T = 0$.
\end{proof}

\begin{exemple}[Champs de Jacobi sur $S^n$]
Sur $S^n$ (courbure $K = 1$), le long d'un grand cercle $\gamma$ param\'etr\'e par
longueur d'arc, l'\'equation de Jacobi pour un champ orthogonal $J \perp \dot\gamma$
devient~:
\[
    J''(t) + J(t) = 0,
\]
de solution $J(t) = A\cos t + B\sin t$. Les points conjugu\'es correspondent
\`a $t = k\pi$, $k \in \Z^*$.
\end{exemple}

\begin{exemple}[Champs de Jacobi sur $\mathbb{H}^n$]
Sur $\mathbb{H}^n$ (courbure $K = -1$), pour $J \perp \dot\gamma$~:
\[
    J''(t) - J(t) = 0,
\]
de solution $J(t) = A\cosh t + B\sinh t$. Il n'y a \emph{pas} de points conjugu\'es.
\end{exemple}

\section{Calcul de la courbure sur les exemples}

\begin{exemple}[Courbure de $S^2$]
En coordonn\'ees sph\'eriques, on calcule~:
\[
    R^1_{212} = \partial_1\Gamma^1_{22} - \partial_2\Gamma^1_{12}
    + \Gamma^1_{1m}\Gamma^m_{22} - \Gamma^1_{2m}\Gamma^m_{12}.
\]
On obtient $R^1_{212} = \sin^2\!\theta$, d'o\`u
$R_{1212} = g_{11} R^1_{212} = \sin^2\!\theta$. La courbure de Gauss est
$K = \frac{R_{1212}}{g_{11}g_{22} - g_{12}^2} = \frac{\sin^2\theta}{\sin^2\theta} = 1$.
\end{exemple}

\begin{exemple}[Courbure de $\mathbb{H}^2$]
Avec la m\'etrique $g = y^{-2}(dx^2 + dy^2)$, on calcule $R_{1212} = -1/y^4$,
$\det(g) = 1/y^4$, d'o\`u $K = -1$.
\end{exemple}

\begin{exemple}[Courbure d'un groupe de Lie bi-invariant]
Pour $X, Y, Z \in \mathfrak{g}$ invariants \`a gauche, avec $\nabla_X Y = \frac{1}{2}[X,Y]$~:
\[
    R(X,Y)Z = \frac{1}{4}[[X,Y],Z].
\]
\end{exemple}

\section{Formule de la seconde variation}

\begin{theoreme}[Seconde variation de l'\'energie]
Soit $\gamma : [0, L] \to M$ une g\'eod\'esique param\'etr\'ee par longueur d'arc et
$\gamma_s$ une variation \`a extr\'emit\'es fixes avec champ de variation $V$.
Alors~:
\[
    \frac{d^2}{ds^2}\bigg|_{s=0} E(\gamma_s) = \int_0^L \left(
    \norm{\frac{DV}{dt}}^2 - \inner{R(\dot\gamma, V)\dot\gamma}{V}\right) dt.
\]
\end{theoreme}

\begin{definition}[Forme d'indice]
La \emph{forme d'indice} de la g\'eod\'esique $\gamma$ est la forme bilin\'eaire~:
\[
    I(V, W) = \int_0^L \left(\inner{\frac{DV}{dt}}{\frac{DW}{dt}}
    - \inner{R(\dot\gamma, V)\dot\gamma}{W}\right) dt.
\]
\end{definition}

\section{Exercices}

\begin{exercice}
V\'erifier la premi\`ere identit\'e de Bianchi $R(X,Y)Z + R(Y,Z)X + R(Z,X)Y = 0$
pour la sph\`ere $S^2$ en coordonn\'ees sph\'eriques.
\end{exercice}

\begin{exercice}
Montrer qu'en dimension $2$, le tenseur de Riemann est enti\`erement d\'etermin\'e
par la courbure de Gauss $K$~:
$R_{1212} = K\,(g_{11}g_{22} - g_{12}^2) = K\det(g)$.
\end{exercice}

\begin{exercice}
Calculer les champs de Jacobi le long d'une g\'eod\'esique radiale du plan
hyperbolique, et v\'erifier qu'il n'y a pas de points conjugu\'es.
\end{exercice}

\begin{exercice}
Pour le groupe de Lie $SO(3)$ avec m\'etrique bi-invariante, calculer $R(X,Y)Z$
pour une base orthonorm\'ee $(X, Y, Z)$ de $\mathfrak{so}(3)$.
\end{exercice}

\begin{exercice}
Montrer que si $(M,g)$ est plate, alors tout champ de Jacobi s'\'ecrit
$J(t) = A + tB$ pour des vecteurs parall\`eles $A$ et $B$.
\end{exercice}

\begin{exercice}
Calculer le tenseur de Riemann du produit tordu $g = dt^2 + f(t)^2 g_{S^{n-1}}$
en fonction de $f$ et de ses d\'eriv\'ees.
\end{exercice}
